% V2 PDF物理页118原书思考题16.1--16.9；图16.18理想电表条件待专家核验。
\section*{第16章\quad 恒定电流：思考题}
\paragraph{思考题 16.1}
导体中没有净电流时能否存在内部电场？导体没有内部电场时是否必然没有电流？区分静电平衡与非稳态情形。

\paragraph{思考题 16.2}
证明同种材料制成、长度相同且截面一致的均匀导线，其电阻与总质量成反比。

\paragraph{思考题 16.3}
半导体和绝缘体的电阻在某些温度范围内随温度升高而下降，怎样从载流子激发和迁移率解释？

\paragraph{思考题 16.4}
从电流密度和电磁能守恒说明 Kirchhoff 第二定律在普通集中参数电路中的成立条件。

\paragraph{思考题 16.5}
电场和电势有什么区别？能否将零电势与零电场混同？

\paragraph{思考题 16.6}
在电阻、电容元件允许的范围内给出一种测量 RC 电路中很高电阻的方案，推导待测电阻公式。

\paragraph{思考题 16.7}
原书图16.18含理想电压表、电流表及5 kΩ、10 kΩ支路，仅凭图形能否唯一估算端点A、B的等效电阻？应说明哪些仪表内阻。

\paragraph{思考题 16.8}
假设人体手至脚约10 kΩ，给出数量级上避免触及未知高电压的物理理由；区分触电电流与电压本身。

\paragraph{思考题 16.9}
分析范德格拉夫静电加速器：传送带怎样将电荷送至金属球？电动势源是哪种非静电作用力？为什么能把机械功转为电能？
